\section{Files d’attente}
Afin de pouvoir avoir une execution de plusieurs threads et de gerer ces executions (thread qui cede le processeur: thread qui termine et plus tard afin de pouvoir appliquer une politique de preemption, 
nous avons ajouter une file d'attente permettant de stocker les threads prets a l'execution. Nous avons opté pour le choix d'une file doublement chainée afin de pourvoir effectué 
des insertions, suppressions et modifications en temps constant. Obtenir une complexité base sur ces operations est important dans une recherche de performance notamment lorsque ces actions sont essentielle
comme par exexmple le retrait d'un thread de la file
lors d'un changement d'etat de ce dernier. Cette operation sur une file simple aurait entrainée une complexité linéaire non négligable.
Ce choix porte un coup, notamment un coup mémoire. Seulement, le nombre de thread restant globalement faible: le surcoup mémoire apporté nous parait negligable.
Notre implémentation de la file repose sur celle décrite par la librarie \textit{<sys/queue.h>}\cite{queue_lib} ous avons ajouté à la structure \textit{thread\_struct} 
